\subsection{全书随堂自测参考解答}
\phantomsection\label{app:chapter-checks}

按正文章号排列，题号与章内自测一致；综合题统一在附录 C 与 D.1。
遇到困难时可先看各章的“可选提示”；参考解答按计算、辨析或证明的需要展开。
标题旁的章号对应正文，标题后注明该章自测所在页；五题数学自检与答案集中在\hyperref[app:math-entry]{数学速查（第~\pageref*{app:math-entry}~页）}。
% Each chapter heading names the page of its self-check box.
\newcommand{\answercheckpage}[1]{\enspace{\mdseries\footnotesize\hyperref[check:#1]{（题目见第~\pageref*{check:#1}~页）}}}

% Answer bodies choose their own calculation, explanation or proof structure.
\newcommand{\checkanswer}[2]{%
  \par\smallskip\noindent
  {\sffamily\bfseries\color{blue}题~#1 \textnormal{·} #2}\par\nobreak
  \noindent\ignorespaces
}
\begingroup
\setlength{\parskip}{3pt}
\begin{enumerate}[leftmargin=5.4em,labelsep=.7em,itemsep=9pt]

\par\Needspace{8\baselineskip}
\item[\textbf{序章}] \phantomsection\label{answer:prologue}\textbf{\hyperref[check:prologue]{序章：为什么 AI 能接住我们说的话？}}\answercheckpage{prologue}\par\nobreak
\begin{answerhints}
\item[\hyperref[check:prologue]{1}] \leavevmode\phantomsection\label{hint:prologue:1}先只看第一个 Token，再每次多看一个。
\item[\hyperref[check:prologue]{2}] \leavevmode\phantomsection\label{hint:prologue:2}对照整句的编号 $[2,4,5,3]$。
\item[\hyperref[check:prologue]{3}] \leavevmode\phantomsection\label{hint:prologue:3}分别问：文字流畅和事实正确由什么检验？
\end{answerhints}
\checkanswer{1}{练习题}共 $\answerresult{3}$ 道：“春天”$\to$“花”；“春天 花”$\to$“开了”；“春天 花 开了”$\to$“。”。最后一个 Token 后面没有答案，所以四个 Token 给出三道题。
\par\checkanswer{2}{编号}“花 开了 。”写成 $\answerresult{[4,5,3]}$；编号 $1$ 留给词表之外的“其他”。
\par\checkanswer{3}{找错}流畅程度是生成形式的表现，既未证明训练资料真实，也未核实新近事实。
\par\Needspace{8\baselineskip}
\item[\textbf{第 1 章}] \phantomsection\label{answer:functions-and-loss}\textbf{\hyperref[check:functions-and-loss]{从可调函数到损失：怎样比较预测？}}\answercheckpage{functions-and-loss}\par\nobreak
\phantomsection\label{answer:loss-analytic}
\begin{answerhints}
\item[\hyperref[check:functions-and-loss]{1}] \leavevmode\phantomsection\label{hint:functions-and-loss:1}先只计算 $w$ 增加的那一份输出。
\item[\hyperref[check:functions-and-loss]{2}] \leavevmode\phantomsection\label{hint:functions-and-loss:2}先问参数有没有被更新。
\item[\hyperref[check:functions-and-loss]{3}] \leavevmode\phantomsection\label{hint:functions-and-loss:3}分别检查前两个样本是否猜对。
\item[\hyperref[check:functions-and-loss]{4}] \leavevmode\phantomsection\label{hint:functions-and-loss:4}按样本顺序计算三项预测。
\item[\hyperref[check:functions-and-loss]{5}] \leavevmode\phantomsection\label{hint:functions-and-loss:5}区分候选集合与全部参数。
\end{answerhints}
\checkanswer{1}{预测}变化量 $(3-2)\cdot3=3$；原输出 $7$，新输出 $\answerresult{10}$。
\checkanswer{2}{找错}推理是固定当前参数计算新输入；训练要用数据、损失和梯度更新参数，不能因更换推理输入便认为已再次训练。
\checkanswer{3}{解释}残差 $1$ 和 $-1$ 相抵，但对应两个错误预测；均方误差为 $(1^2+(-1)^2+0^2)/3=\answerresult{2/3}>0$。
\checkanswer{4}{手算与辨析}$w=2,b=1$ 的预测为 $\answerresult{(-1,1,3)}$，残差为 $(0,-1,1)$，均方误差 $\answerresult{2/3}$，与候选 B 相同。本题用相同损失比较不同预测：B 的预测是 $(0,1,2)$，误差分布与这条直线不同。两组参数在损失等高线上有相同高度，并未得到同一个函数；在新输入 $x=2$ 处，两者还分别预测 $3$ 与 $5$。
\checkanswer{5}{找错}B 只是已经比较的三个候选中损失最小的一个。要说明全局最优，还需证明其他任意参数都无法得到更小的损失。
\par\Needspace{8\baselineskip}
\item[\textbf{第 2 章}] \phantomsection\label{answer:gradient-descent}\textbf{\hyperref[check:gradient-descent]{怎样让损失变小？梯度下降调参数}}\answercheckpage{gradient-descent}\par\nobreak
 \begin{answerhints}
\item[\hyperref[check:gradient-descent]{1}] \leavevmode\phantomsection\label{hint:gradient-descent:1}把学习率零代入更新式。
\item[\hyperref[check:gradient-descent]{2}] \leavevmode\phantomsection\label{hint:gradient-descent:2}先在旧参数处同时求两个偏导。
\item[\hyperref[check:gradient-descent]{3}] \leavevmode\phantomsection\label{hint:gradient-descent:3}圈出本次梯度使用的是哪一版参数。
\end{answerhints}
\checkanswer{1}{预测}学习率为零时仍处于 $(1,0)$，损失为 $5/3$；梯度 $(-2/3,-2)$ 非零，参数不动是因为步长为零，并非到达最优点。
\par\checkanswer{2}{手算}$\nabla L=(-2/3,-2)$；以 $\eta=0.3$ 同步更新到 $(1.2,0.6)$。新预测为 $(-0.6,0.6,1.8)$，残差为 $(0.4,-1.4,-0.2)$，损失为 $(0.4^2+(-1.4)^2+(-0.2)^2)/3=\answerresult{0.72}$。
\par\checkanswer{3}{找错}同步更新要在同一旧参数处算完全部梯度后再一起更新；不能将新参数混进本次剩余梯度。
\par\Needspace{8\baselineskip}
\item[\textbf{第 3 章}] \phantomsection\label{answer:nonlinearity}\textbf{\hyperref[check:nonlinearity]{非线性：为什么多叠两层直线仍不够？}}\answercheckpage{nonlinearity}\par\nobreak
\begin{answerhints}
\item[\hyperref[check:nonlinearity]{1}] \leavevmode\phantomsection\label{hint:nonlinearity:1}先展开两个仿射变换的复合式。
\item[\hyperref[check:nonlinearity]{2}] \leavevmode\phantomsection\label{hint:nonlinearity:2}先算 ReLU 括号里的数。
\item[\hyperref[check:nonlinearity]{3}] \leavevmode\phantomsection\label{hint:nonlinearity:3}分别看折角两侧的斜率。
\end{answerhints}
\checkanswer{1}{比较}$(ax+b)c+d=(ac)x+(bc+d)$，仍为一次仿射变换。
\par\checkanswer{2}{手算}$f(0)=\relu(-1)=\answerresult{0}$；$f(1)=\relu(1)=\answerresult{1}$。
\par\checkanswer{3}{预测}左侧为常值零、斜率 $0$；右侧斜率为 $2$，因此可以表示折角两边不同的变化率。
\par\Needspace{8\baselineskip}
\item[\textbf{第 4 章}] \phantomsection\label{answer:backpropagation}\textbf{\hyperref[check:backpropagation]{反向传播究竟在传什么？}}\answercheckpage{backpropagation}\par\nobreak
\begin{answerhints}
\item[\hyperref[check:backpropagation]{1}] \leavevmode\phantomsection\label{hint:backpropagation:1}沿 $w_1\to h\to\widehat{y}$ 逐段求导。
\item[\hyperref[check:backpropagation]{2}] \leavevmode\phantomsection\label{hint:backpropagation:2}先写出负区间 ReLU 的局部导数。
\item[\hyperref[check:backpropagation]{3}] \leavevmode\phantomsection\label{hint:backpropagation:3}把从损失到参数的路径画全。
\item[\hyperref[check:backpropagation]{4}] \leavevmode\phantomsection\label{hint:backpropagation:4}按前向、损失、反传、更新排列动作。
\item[\hyperref[check:backpropagation]{5}] \leavevmode\phantomsection\label{hint:backpropagation:5}先求隐藏激活，再代入输出层。
\item[\hyperref[check:backpropagation]{6}] \leavevmode\phantomsection\label{hint:backpropagation:6}先算 $z=w_1x+b_1$ 并判断正负。
\end{answerhints}
\checkanswer{1}{计算}$\partial\widehat{y}/\partial w_1=w_2x=\answerresult{6}$。
\par\checkanswer{2}{预测}$z<0$ 时 ReLU 的局部导数为零，所以该分支反传梯度为 $\answerresult{0}$。
\par\checkanswer{3}{找错}漏掉 $\partial h/\partial w_1$；应求 $(\partial L/\partial h)(\partial h/\partial w_1)$。
\checkanswer{4}{流程}前向产生预测，损失读取标签，反传产生各参数梯度，更新才改变参数。
\par\checkanswer{5}{手算}$h=2,\widehat{y}=2\cdot2=\answerresult{4}$；损失为 $(4-5)^2/2=\answerresult{0.5}$。
\par\checkanswer{6}{手算}$z=1\cdot(-1)+0=\answerresult{-1}$，$h=\relu(-1)=0$，$\widehat{y}=2\cdot0+0=\answerresult{0}$，$\partial L/\partial\widehat{y}=0-5=\answerresult{-5}$。$\partial L/\partial b_2=-5$ 仍非零；$\partial L/\partial w_2=-5\cdot h=0$；ReLU 在 $z<0$ 处导数为零，所以 $\partial L/\partial w_1=\partial L/\partial b_1=0$。
\par\Needspace{8\baselineskip}
\item[\textbf{第 5 章}] \phantomsection\label{answer:generalization}\textbf{\hyperref[check:generalization]{训练得很好，为什么新题反而做错？}}\answercheckpage{generalization}\par\nobreak
\begin{answerhints}
\item[\hyperref[check:generalization]{1}] \leavevmode\phantomsection\label{hint:generalization:1}对比训练与验证曲线的方向。
\item[\hyperref[check:generalization]{2}] \leavevmode\phantomsection\label{hint:generalization:2}逐轮比较验证损失。
\item[\hyperref[check:generalization]{3}] \leavevmode\phantomsection\label{hint:generalization:3}先查 Dropout 屏蔽的对象及测试集用途。
\end{answerhints}
\checkanswer{1}{预测}训练降、验证升提示先排查过拟合，应先检查划分与泄漏，再考虑早停、正则化等。
\par\checkanswer{2}{手算}第 $\answerresult{2}$ 轮验证损失最低，为 $\answerresult{0.4}$。
\par\checkanswer{3}{找错}Dropout 训练时随机屏蔽\textbf{激活}，不是永久删除参数；超参数在验证数据上选择，不可反复用最终测试集挑选。
\par\Needspace{8\baselineskip}
\item[\textbf{第 6 章}] \phantomsection\label{answer:vectors}\textbf{\hyperref[check:vectors]{向量与点积：怎样把多个特征变成一个匹配分数？}}\answercheckpage{vectors}\par\nobreak
\begin{answerhints}
\item[\hyperref[check:vectors]{1}] \leavevmode\phantomsection\label{hint:vectors:1}先为两个位置写上“技术、生活”。
\item[\hyperref[check:vectors]{2}] \leavevmode\phantomsection\label{hint:vectors:2}将需求与新文章的两个分量分别相乘。
\item[\hyperref[check:vectors]{3}] \leavevmode\phantomsection\label{hint:vectors:3}先比较放大前后的点积，再看方向。
\item[\hyperref[check:vectors]{4}] \leavevmode\phantomsection\label{hint:vectors:4}先用勾股定理算长度。
\item[\hyperref[check:vectors]{5}] \leavevmode\phantomsection\label{hint:vectors:5}由单位向量点积等于夹角余弦判断。
\item[\hyperref[check:vectors]{6}] \leavevmode\phantomsection\label{hint:vectors:6}先回看长度的作用，再问表示怎样得到。
\end{answerhints}
\checkanswer{1}{表示}原需求的生活分数 $0.8$ 大于技术分数 $0.6$；交换后技术分数更高，变成\textbf{稍偏向技术}。只换数字而保持特征顺序，就改变了需求。
\par\checkanswer{2}{手算}$0.6\times0.6+0.8\times0.8=0.36+0.64=\answerresult{1}$，超过 A 的 $0.96$，因此新文章排在 A 前。两条向量相同且长度均为 $1$，这个分数也对应方向完全一致。
\par\checkanswer{3}{预测}原分数 $0.6$，放大后为 $0.6\times3+0.8\times0=\answerresult{1.8}$，增至三倍；两个分量仍只有技术项非零，\textbf{主题配比不变}。
\par\checkanswer{4}{手算}$\sqrt{3^2+0^2}=\answerresult{3}$；逐项除以 $3$ 得到$\answerresult{(1,0)}$。需求本身长为 $1$，所以归一化后点积为 $0.6\times1+0.8\times0=\answerresult{0.6}$，也等于原始点积 $1.8$ 除以长度乘积 $1\times3$。
\par\checkanswer{5}{几何}点积为零说明夹角为直角，两条向量\textbf{垂直}；其中一条反向后仍然垂直，新点积仍为$\answerresult{0}$。
\par\checkanswer{6}{找错}首先核对\textbf{长度是否影响了比较}，原始点积会随整体放大而增大；其次核对\textbf{表示是否适合所需的语义比较}，同维、同尺度的数字仍须通过合适的训练或特征设计反映任务关系。
\par\Needspace{8\baselineskip}
\item[\textbf{第 7 章}] \phantomsection\label{answer:matrix-multiplication}\textbf{\hyperref[check:matrix-multiplication]{矩阵乘法到底在算什么？}}\answercheckpage{matrix-multiplication}\par\nobreak
\begin{answerhints}
\item[\hyperref[check:matrix-multiplication]{1}] \leavevmode\phantomsection\label{hint:matrix-multiplication:1}先比较内侧两个维度。
\item[\hyperref[check:matrix-multiplication]{2}] \leavevmode\phantomsection\label{hint:matrix-multiplication:2}将 $(1,0)$ 分别与三篇文章的特征做点积。
\item[\hyperref[check:matrix-multiplication]{3}] \leavevmode\phantomsection\label{hint:matrix-multiplication:3}先检查交换后内侧维度是否相等。
\end{answerhints}
\checkanswer{1}{形状}$2\times3$ 乘 $3\times4$ 内侧维度相同，输出 $\answerresult{2\times4}$；若右矩阵变为 $4\times3$，内侧 $3\ne4$，不合法。
\par\checkanswer{2}{手算}得到 $\answerresult{(0.8,1,0)}$，对应 A、B、C。文章按行放在 $K:3\times2$ 中，转置后每列一篇，才能与需求行逐项配对；$(1\times2)(2\times3)\to1\times3$。
\par\checkanswer{3}{找错}矩阵乘法不满足一般交换律。前例 $AB$ 合法，但 $BA$ 连形状也不合法；即使均合法，可取具体数字检验结果是否相等。
\par\Needspace{8\baselineskip}
\item[\textbf{第 8 章}] \phantomsection\label{answer:mlp}\textbf{\hyperref[check:mlp]{MLP：把矩阵与激活函数组成两层网络}}\answercheckpage{mlp}\par\nobreak
\begin{answerhints}
\item[\hyperref[check:mlp]{1}] \leavevmode\phantomsection\label{hint:mlp:1}先算 $xW_1$ 的形状。
\item[\hyperref[check:mlp]{2}] \leavevmode\phantomsection\label{hint:mlp:2}先让 $(1,-2)$ 分别点乘 $W_1$ 的三列。
\item[\hyperref[check:mlp]{3}] \leavevmode\phantomsection\label{hint:mlp:3}先检查公式中有没有读取别的行。
\end{answerhints}
\checkanswer{1}{形状}$xW_1:1\times3$，随后 $xW_1W_2:1\times2$。
\par\checkanswer{2}{手算}$\vect z^\top=(1,\,-2,\,1-2)=\answerresult{(1,-2,-1)}$；ReLU 后 $\vect h^\top=\answerresult{(1,0,0)}$，第 2、3 个隐藏单元被截断为零；$\widehat{\vect y}^\top=\vect h^\top W_2=\answerresult{(1,0)}$，只剩 $W_2$ 第一行起作用。
\par\checkanswer{3}{找错}参数共享让各行使用同一函数；逐 Token FFN 对每行单独计算，Attention 等机制才负责跨位置读取。
\par\Needspace{8\baselineskip}
\item[\textbf{第 9 章}] \phantomsection\label{answer:tokenization}\textbf{\hyperref[check:tokenization]{Token 与 ID：文字怎样进入数字世界？}}\answercheckpage{tokenization}\par\nobreak
\begin{answerhints}
\item[\hyperref[check:tokenization]{1}] \leavevmode\phantomsection\label{hint:tokenization:1}先把 ID 看作词表标签。
\item[\hyperref[check:tokenization]{2}] \leavevmode\phantomsection\label{hint:tokenization:2}按给定合并规则从左到右切分。
\item[\hyperref[check:tokenization]{3}] \leavevmode\phantomsection\label{hint:tokenization:3}先检查切分结果依赖哪些输入。
\end{answerhints}
\checkanswer{1}{解释}ID 是查表用的离散编号，不是可比较大小的语义强度。
\par\checkanswer{2}{手算}按指定合并得到“今天／天气／好”，共 $\answerresult{3}$ 个 Token。
\par\checkanswer{3}{找错}BPE 合并规则与词表由训练确定，实际切分还取决于新文本；不同文本或词表不保证相同 Token 数。
\par\Needspace{8\baselineskip}
\item[\textbf{第~\ref*{intro:language-model-task}~章}] \phantomsection\label{answer:language-model-task}\textbf{\hyperref[check:language-model-task]{下一 Token 预测：训练与生成为何不同？}}\answercheckpage{language-model-task}\par\nobreak
\begin{answerhints}
\item[\hyperref[check:language-model-task]{1}] \leavevmode\phantomsection\label{hint:language-model-task:1}在完整文本中找到每个输入 Token，取它右边紧邻的一项作目标。
\item[\hyperref[check:language-model-task]{2}] \leavevmode\phantomsection\label{hint:language-model-task:2}先数预测位置，再数每行候选。
\item[\hyperref[check:language-model-task]{3}] \leavevmode\phantomsection\label{hint:language-model-task:3}先写出第一个位置的目标，再检查它是否已出现在预测所用的输入中。
\end{answerhints}
\checkanswer{1}{对齐}位置~1 的“春天”（ID $2$）后面是“花”（ID $4$）；位置~2 的“花”（ID $4$）后面是“开了”（ID $5$）；位置~3 的“开了”（ID $5$）后面是句号（ID $3$）。按预测位置排列，三个目标是 $\answerresult{[4,5,3]}$。
\par\checkanswer{2}{形状}三个位置乘五类分数，分数矩阵为 $\answerresult{3\times5}$。
\par\checkanswer{3}{找错}第一个位置的目标就是“花”，本应只利用“春天”来预测；读取第二位置的“花”等于提前看到了答案。目标在评分步骤参与损失计算。
\par\Needspace{8\baselineskip}
\item[\textbf{第~\ref*{intro:classification-loss}~章}] \phantomsection\label{answer:classification-loss}\textbf{\hyperref[check:classification-loss]{分类损失：预测错了哪一类，怎样调整分数？}}\answercheckpage{classification-loss}\par\nobreak
\begin{answerhints}
\item[\hyperref[check:classification-loss]{1}] \leavevmode\phantomsection\label{hint:classification-loss:1}区分选出一个候选，与衡量正确答案获得多少支持。
\item[\hyperref[check:classification-loss]{2}] \leavevmode\phantomsection\label{hint:classification-loss:2}逐项检查非负性，再求和。
\item[\hyperref[check:classification-loss]{3}] \leavevmode\phantomsection\label{hint:classification-loss:3}先找到正确类别的那一项概率。
\item[\hyperref[check:classification-loss]{4}] \leavevmode\phantomsection\label{hint:classification-loss:4}比较 $-\ln p$ 随 $p$ 增大的方向。
\end{answerhints}
\checkanswer{1}{辨析}只选最高分时，直接比较 $(2,-1,0)$ 就能选第一类。本章训练时采用概率加交叉熵的方案：先把 logits 变成总和为 $1$ 的概率，再按正确答案获得的概率计算损失。这样，即使两次都选对，也能区分正确答案概率为 $0.4$ 与 $0.9$ 的预测，后者损失更小。
\par\checkanswer{2}{判断}只有 $\answerresult{(0.2,0.5,0.3)}$ 合法：三项非负且总和为 $1$。第二组总和为 $1.1$；第三组虽总和为 $1$，却含负数，都不能作为类别概率。
\par\checkanswer{3}{手算}$-\ln(1/2)=\answerresult{\ln2\approx0.6931}$，该类梯度 $1/2-1=\answerresult{-0.5}$。
\par\checkanswer{4}{预测}$-\ln0.4<-\ln0.2$，交叉熵下降。
\par\Needspace{8\baselineskip}
\item[\textbf{第~\ref*{intro:embeddings}~章}] \phantomsection\label{answer:embeddings}\textbf{\hyperref[check:embeddings]{Embedding：让 Token ID 变成可训练的向量}}\answercheckpage{embeddings}\par\nobreak
\begin{answerhints}
\item[\hyperref[check:embeddings]{1}] \leavevmode\phantomsection\label{hint:embeddings:1}先数查出的行数与每行宽度。
\item[\hyperref[check:embeddings]{2}] \leavevmode\phantomsection\label{hint:embeddings:2}把 one-hot 与表的每一行分别配对。
\item[\hyperref[check:embeddings]{3}] \leavevmode\phantomsection\label{hint:embeddings:3}先找到两个位置读取的表行。
\end{answerhints}
\checkanswer{1}{形状}Embedding 表为 $5\times3$，三个 ID 查表为 $\answerresult{3\times3}$，每行属于输入的一个位置。
\par\checkanswer{2}{比较}ID 只是任意编号；one-hot 行仅有对应 ID 列为 $1$，右乘 Embedding 表正好选出对应行，与查表同义。
\par\checkanswer{3}{预测}重复 ID 的原始查表向量相同；训练更新这行后，下次两个位置都读到同一条更新后的向量。位置不同本身不会改变查表结果。
\par\Needspace{8\baselineskip}
\item[\textbf{第 13 章}] \phantomsection\label{answer:context-and-rnn}\textbf{\hyperref[check:context-and-rnn]{前文怎样传到当前词？从逐行处理到 RNN}}\answercheckpage{context-and-rnn}\par\nobreak
\begin{answerhints}
\item[\hyperref[check:context-and-rnn]{1}] \leavevmode\phantomsection\label{hint:context-and-rnn:1}先比较两句传给当前 Token 函数的输入。
\item[\hyperref[check:context-and-rnn]{2}] \leavevmode\phantomsection\label{hint:context-and-rnn:2}先把 $x=3$ 代入同一个函数。
\item[\hyperref[check:context-and-rnn]{3}] \leavevmode\phantomsection\label{hint:context-and-rnn:3}先问可变历史如何放进固定宽度输入。
\item[\hyperref[check:context-and-rnn]{4}] \leavevmode\phantomsection\label{hint:context-and-rnn:4}先看公式里前一步的状态从哪里进入。
\item[\hyperref[check:context-and-rnn]{5}] \leavevmode\phantomsection\label{hint:context-and-rnn:5}从 $h_0=0$ 起逐步代入。
\item[\hyperref[check:context-and-rnn]{6}] \leavevmode\phantomsection\label{hint:context-and-rnn:6}检查更新状态所需的先后依赖。
\end{answerhints}
\checkanswer{1}{预测}不能。输入只保留当前 Token 时，两句进入函数的是同一个表示，输出无法仅因不可见的前文改变。
\par\checkanswer{2}{手算}两句都得到 $\answerresult{7}$。
\par\checkanswer{3}{找错}需要固定输入长度／截断和相应参数或计算预算；直接拼接不能免费覆盖无限历史。
\par\checkanswer{4}{流程}每一步根据 $x_t$ 与 $h_{t-1}$ 计算 $h_t$；状态把部分历史摘要传给未来位置。
\par\checkanswer{5}{手算}$h_1=\answerresult{2}$，$h_2=0+0.5\cdot2=\answerresult{1}$，$h_3=1+0.5\cdot1=\answerresult{1.5}$。
\par\checkanswer{6}{预测}第一个输入增加 $2$，经两次系数 $0.5$ 传播，到 $h_3$ 增加 $2(0.5)^2=\answerresult{0.5}$，变为 $2$。
\par\Needspace{8\baselineskip}
\item[\textbf{第 14 章}] \phantomsection\label{answer:attention-qkv}\textbf{\hyperref[check:attention-qkv]{Attention：从直接读取到 Q/K/V 分工}}\answercheckpage{attention-qkv}\par\nobreak
\phantomsection\label{answer:why-attention}
\begin{answerhints}
\item[\hyperref[check:attention-qkv]{1}] \leavevmode\phantomsection\label{hint:attention-qkv:1}先用第一行和各行做点积。
\item[\hyperref[check:attention-qkv]{2}] \leavevmode\phantomsection\label{hint:attention-qkv:2}分别写出两行 Softmax 的分母。
\item[\hyperref[check:attention-qkv]{3}] \leavevmode\phantomsection\label{hint:attention-qkv:3}先问 Query 和 Key 是否有独立投影。
\item[\hyperref[check:attention-qkv]{4}] \leavevmode\phantomsection\label{hint:attention-qkv:4}先给 Q、K、V 分配查询、匹配、内容角色。
\item[\hyperref[check:attention-qkv]{5}] \leavevmode\phantomsection\label{hint:attention-qkv:5}将两个 Value 分别乘对应权重。
\item[\hyperref[check:attention-qkv]{6}] \leavevmode\phantomsection\label{hint:attention-qkv:6}先找 Value 出现在分数公式还是读出公式中。
\end{answerhints}
\checkanswer{1}{手算}第一行 $(1,0,1)$，第三行 $(1,1,2)$，$S_{13}=S_{31}=\answerresult{1}$。Softmax 保持同一行的大小关系：第一行在位置 $1,3$ 并列最大；第二行 $(0,1,1)$ 在位置 $2,3$ 并列最大；第三行只有位置 $3$ 最大。
\checkanswer{2}{预测}不必相等。$A_{13}=e/(2e+1)$，$A_{31}=e/(2e+e^2)$；虽然对应原始分数同为 $1$，每行 Softmax 的分母不同。
\checkanswer{3}{找错}零参数 $XX^\top$ 让 Query/Key 使用同样的输入向量，不能单独学“如何问”和“如何匹配”。
\checkanswer{4}{定位}Query 发问、Key 用于匹配、Value 提供读取内容；$\vect q^\top K^\top:1\times3$，乘 $V$ 后输出 $1\times2$。
\checkanswer{5}{手算}$0.2(1,0)+0.8(0,1)=\answerresult{(0.2,0.8)}$。
\checkanswer{6}{找错}分数来自 Query 与 Key，Value 只进加权求和。改 Value 不直接改 Q/K 原始分数，但可能改变读出。
\par\Needspace{8\baselineskip}
\item[\textbf{第 15 章}] \phantomsection\label{answer:self-attention}\textbf{\hyperref[check:self-attention]{从单 Query 到批量 Self-Attention}}\answercheckpage{self-attention}\par\nobreak
\begin{answerhints}
\item[\hyperref[check:self-attention]{1}] \leavevmode\phantomsection\label{hint:self-attention:1}先算 $QK^\top$，再接上 $V$。
\item[\hyperref[check:self-attention]{2}] \leavevmode\phantomsection\label{hint:self-attention:2}先比较相同分数的两个指数。
\item[\hyperref[check:self-attention]{3}] \leavevmode\phantomsection\label{hint:self-attention:3}分别试只改一项和同时平移两项。
\end{answerhints}
\checkanswer{1}{形状}匹配分数及 Softmax 均为 $3\times4$，乘 $V:4\times3$ 后输出 $\answerresult{3\times3}$。
\par\checkanswer{2}{手算}两个相同分数对应权重 $(1/2,1/2)$，读出为 $\answerresult{(0.5,0.5)}$。
\par\checkanswer{3}{预测}第一权重增大；若同时加相同常数，分子分母共同缩放，权重不变。
\par\Needspace{8\baselineskip}
\item[\textbf{第 16 章}] \phantomsection\label{answer:position-multihead}\textbf{\hyperref[check:position-multihead]{位置编码、缩放与多头注意力的手算}}\answercheckpage{position-multihead}\par\nobreak
\begin{answerhints}
\item[\hyperref[check:position-multihead]{1}] \leavevmode\phantomsection\label{hint:position-multihead:1}头 2 的缩放分数变成 $(0,1)/\sqrt2$，与头 1 相同。
\item[\hyperref[check:position-multihead]{2}] \leavevmode\phantomsection\label{hint:position-multihead:2}先把两个头的特征列拼起来。
\item[\hyperref[check:position-multihead]{3}] \leavevmode\phantomsection\label{hint:position-multihead:3}给原始分数取一个很大的正数试算。
\end{answerhints}
\checkanswer{1}{手算}权重为 $\answerresult{(0.330,0.670)}$，读出 $0.330(0,1)+0.670(1,0)\approx\answerresult{(0.670,0.330)}$。两个头的最大权重都落在位置 2，只是 Value 不同；各头有独立的 Query 投影，才能像正文那样重点读取不同位置。
\par\checkanswer{2}{形状}拼接为 $\answerresult{3\times4}$；输出投影后为 $\answerresult{3\times6}$。
\par\checkanswer{3}{找错}除以 $\sqrt{d_k}$ 只是缩放，原分数没有保证上界；只有 Softmax 之后的概率落在 $[0,1]$。
\par\Needspace{8\baselineskip}
\item[\textbf{第 17 章}] \phantomsection\label{answer:transformer-architecture}\textbf{\hyperref[check:transformer-architecture]{把 Attention 组装起来：Transformer 不只是一个公式}}\answercheckpage{transformer-architecture}\par\nobreak
\begin{answerhints}
\item[\hyperref[check:transformer-architecture]{1}] \leavevmode\phantomsection\label{hint:transformer-architecture:1}先数拼接后的宽度，再核对残差两端。
\item[\hyperref[check:transformer-architecture]{2}] \leavevmode\phantomsection\label{hint:transformer-architecture:2}在位置图上标出被遮住的上三角区域。
\item[\hyperref[check:transformer-architecture]{3}] \leavevmode\phantomsection\label{hint:transformer-architecture:3}先画出跳过子层的输入旁路。
\end{answerhints}
\checkanswer{1}{形状}拼接得 $\answerresult{3\times4}$；残差相加要求两侧形状兼容，否则需投影对齐。
\par\checkanswer{2}{预测}训练时位置 $2$ 的输入正是位置 $1$ 的目标；位置 $1$ 若能读取它，就能直接复制答案来降低损失，“只用当前及之前 Token 预测下一 Token”的规则随之失效。生成时没有这份未来输入可读，训练所学也就无法用于生成。
\par\checkanswer{3}{找错}本书采用输入 $\to$ Norm $\to$ 子层计算，再把该子层的原输入加回来，即 $x+f(\operatorname{Norm}(x))$。相加后才 Norm 的排列属于 Post-Norm。
\par\Needspace{8\baselineskip}
\item[\textbf{第 18 章}] \phantomsection\label{answer:ffn-output-head}\textbf{\hyperref[check:ffn-output-head]{前馈网络与输出头：两种不同的特征变换}}\answercheckpage{ffn-output-head}\par\nobreak
\begin{answerhints}
\item[\hyperref[check:ffn-output-head]{1}] \leavevmode\phantomsection\label{hint:ffn-output-head:1}先看两种输出各自属于隐藏空间还是词表。
\item[\hyperref[check:ffn-output-head]{2}] \leavevmode\phantomsection\label{hint:ffn-output-head:2}先写出上投影和下投影的内侧维度。
\item[\hyperref[check:ffn-output-head]{3}] \leavevmode\phantomsection\label{hint:ffn-output-head:3}先比较隐藏宽度和词表大小。
\end{answerhints}
\checkanswer{1}{定位}FFN 在模型隐藏空间内映射特征；词表输出头把隐藏向量变为各 Token 类别的 logits。
\par\checkanswer{2}{手算}$W_{\rm up}:6\times18$、$W_{\rm down}:18\times6$，整体输出 $\answerresult{4\times6}$。
\par\checkanswer{3}{找错}FFN 后仍是隐藏向量，还需输出头投影到词表大小，再经 Softmax 得到概率。
\par\Needspace{8\baselineskip}
\item[\textbf{第 19 章}] \phantomsection\label{answer:transformer-forward-pass}\textbf{\hyperref[check:transformer-forward-pass]{三维贯通实验：从 Token ID 到三个位置的输出概率}}\answercheckpage{transformer-forward-pass}\par\nobreak
 \begin{answerhints}
\item[\hyperref[check:transformer-forward-pass]{1}] \leavevmode\phantomsection\label{hint:transformer-forward-pass:1}沿形状地图找出首次从三列变成五列的位置。
\item[\hyperref[check:transformer-forward-pass]{2}] \leavevmode\phantomsection\label{hint:transformer-forward-pass:2}先把最高概率 ID 与目标 ID 逐行对照。
\item[\hyperref[check:transformer-forward-pass]{3}] \leavevmode\phantomsection\label{hint:transformer-forward-pass:3}先在公式里找跨行读取发生在哪一步。
\end{answerhints}
\checkanswer{1}{定位}$X,A,H:3\times3$，$Z,\Pi:3\times5$；输出头 $W_{\rm out}:3\times5$ 从隐藏特征生成词表分数。
\par\checkanswer{2}{预测}初始位置 $1,2$ 错，位置 $3$ 对。因果 Attention 允许第三位置读取第一位置，所以首 Token 改变可影响末位置 logits。
\par\checkanswer{3}{找错}说反了：Attention 跨 Token 位置读取；FFN 使用共享权重逐行做非线性变换，不交换 Token 信息。
\par\Needspace{8\baselineskip}
\item[\textbf{第 20 章}] \phantomsection\label{answer:language-model-training}\textbf{\hyperref[check:language-model-training]{三个位置一起训练：从交叉熵到共享参数更新}}\answercheckpage{language-model-training}\par\nobreak
\begin{answerhints}
\item[\hyperref[check:language-model-training]{1}] \leavevmode\phantomsection\label{hint:language-model-training:1}将输入与右移一位的目标逐行配对。
\item[\hyperref[check:language-model-training]{2}] \leavevmode\phantomsection\label{hint:language-model-training:2}先写出 $(p-y)/N$。
\item[\hyperref[check:language-model-training]{3}] \leavevmode\phantomsection\label{hint:language-model-training:3}先给每份参数标上更新步号。
\end{answerhints}
\checkanswer{1}{对齐}三位置目标为 $\answerresult{[4,5,3]}$；因果掩码 限制 Attention 可见位置，损失掩码 选择参与计分的位置。
\par\checkanswer{2}{手算}$(0.3-1)/3=\answerresult{-0.2333}$；若仅两位置参与，$(0.3-1)/2=\answerresult{-0.35}$。
\par\checkanswer{3}{找错}同一训练步必须在一套旧参数上算全部梯度，再一起更新；混用更新后的输出头计算前层梯度不属于同一次同步更新。
\par\Needspace{8\baselineskip}
\item[\textbf{第 21 章}] \phantomsection\label{answer:autoregressive-generation}\textbf{\hyperref[check:autoregressive-generation]{训练好了，模型怎样逐词生成文字？}}\answercheckpage{autoregressive-generation}\par\nobreak
 \begin{answerhints}
\item[\hyperref[check:autoregressive-generation]{1}] \leavevmode\phantomsection\label{hint:autoregressive-generation:1}先看除以更小温度后分数差如何变化。
\item[\hyperref[check:autoregressive-generation]{2}] \leavevmode\phantomsection\label{hint:autoregressive-generation:2}按概率降序累加到阈值，再重新归一化。
\item[\hyperref[check:autoregressive-generation]{3}] \leavevmode\phantomsection\label{hint:autoregressive-generation:3}先区分已知训练前缀与尚未生成的新 Token。
\end{answerhints}
\checkanswer{1}{预测}C 从 $\tau=1$ 时约 $0.090$ 降至 $\tau=0.5$ 时约 $\answerresult{0.016}$；$\tau$ 趋于无穷时三个概率均趋 $1/3$。
\par\checkanswer{2}{手算}Top-$p=0.8$ 累计 $0.5+0.3=0.8$，按大于等于阈值保留 A、B，归一化为 $\answerresult{(0.625,0.375)}$。Top-$k=3$ 保留前三，总和 $0.95$，归一化约 $\answerresult{(0.526,0.316,0.158)}$；保留集合不同。
\par\checkanswer{3}{找错}训练时真实前缀全部已知，因果掩码允许不同位置同一轮前向预测；生成时新输入由上一步选出的未知 Token 决定，必须顺序生成。
\par\Needspace{8\baselineskip}
\item[\textbf{第 22 章}] \phantomsection\label{answer:real-models}\textbf{\hyperref[check:real-models]{从三维到真实模型：规模、配置与分词}}\answercheckpage{real-models}\par\nobreak
\begin{answerhints}
\item[\hyperref[check:real-models]{1}] \leavevmode\phantomsection\label{hint:real-models:1}用模型宽度除以头数。
\item[\hyperref[check:real-models]{2}] \leavevmode\phantomsection\label{hint:real-models:2}分别写出四张 Attention 投影和两张 FFN 的元素数。
\item[\hyperref[check:real-models]{3}] \leavevmode\phantomsection\label{hint:real-models:3}先问损失来自哪份数据和哪种目标。
\end{answerhints}
\checkanswer{1}{形状}$768/12=\answerresult{64}$；每头宽度不直接给出层数。
\par\checkanswer{2}{估算}Attention 约 $4\cdot768^2$，FFN 约 $8\cdot768^2$，后者约为前者两倍（按题设简化）。
\par\checkanswer{3}{找错}训练损失仅反映对应数据目标，还需独立泛化评估与事实／鲁棒性等任务检查。
\par\Needspace{8\baselineskip}
\item[\textbf{第 23 章}] \phantomsection\label{answer:pretraining-icl}\textbf{\hyperref[check:pretraining-icl]{预训练已经学会续写，为什么还能临时学会新任务？}}\answercheckpage{pretraining-icl}\par\nobreak
\begin{answerhints}
\item[\hyperref[check:pretraining-icl]{1}] \leavevmode\phantomsection\label{hint:pretraining-icl:1}逐项检查是否发生损失反传和参数更新。
\item[\hyperref[check:pretraining-icl]{2}] \leavevmode\phantomsection\label{hint:pretraining-icl:2}先数当前 Prompt 里的任务记录。
\item[\hyperref[check:pretraining-icl]{3}] \leavevmode\phantomsection\label{hint:pretraining-icl:3}先检查复制示例后是否执行了优化步骤。
\end{answerhints}
\checkanswer{1}{比较}预训练通过损失反传更新参数；本章的 Few-shot ICL 将样例加进当前上下文，权重保持固定。
\par\checkanswer{2}{手算}两例加一待答共 $\answerresult{3}$ 条记录；另一会话没有带入时，原例不在新的上下文中。
\par\checkanswer{3}{找错}Few-shot 不永久改写权重；新 Q3 数据须来自有权限、可核验的报告或工具读取。
\par\Needspace{8\baselineskip}
\item[\textbf{第 24 章}] \phantomsection\label{answer:post-training}\textbf{\hyperref[check:post-training]{怎样改进回答：训练反馈与推理时计算}}\answercheckpage{post-training}\par\nobreak
\phantomsection\label{answer:test-time-compute}
\begin{answerhints}
\item[\hyperref[check:post-training]{1}] \leavevmode\phantomsection\label{hint:post-training:1}先检查是否反传更新，再分别追踪三类反馈来源。
\item[\hyperref[check:post-training]{2}] \leavevmode\phantomsection\label{hint:post-training:2}先算交叉熵，再求奖励均值与标准差。
\item[\hyperref[check:post-training]{3}] \leavevmode\phantomsection\label{hint:post-training:3}先查相对优势之外还有哪些策略目标项。
\item[\hyperref[check:post-training]{4}] \leavevmode\phantomsection\label{hint:post-training:4}先检查参数有没有被更新。
\item[\hyperref[check:post-training]{5}] \leavevmode\phantomsection\label{hint:post-training:5}先计算四次都失败的概率。
\item[\hyperref[check:post-training]{6}] \leavevmode\phantomsection\label{hint:post-training:6}枚举四个独立二值结果，再对比一个共享结果。
\end{answerhints}
\checkanswer{1}{比较}预训练与 SFT 更新参数；Few-shot ICL 只改变当前上下文，参数固定。SFT 使用示范答案；偏好学习使用回答比较；可验证奖励使用程序、判题器等可复现的核验结果。
\checkanswer{2}{手算}$-\ln0.25=\answerresult{\ln4\approx1.3863}$；奖励均值 $0.5$、总体标准差 $0.5$，优势为 $\answerresult{(1,-1,1,-1)}$。
\checkanswer{3}{找错}完整目标还涉及新旧策略概率比、裁剪和可能的正则化；奖励完全相同时加稳定项可避免除零，但所有优势仍为零，这组样本的策略项不提供区分性梯度。
\checkanswer{4}{解释}RLVR 用验证结果构造奖励并更新策略；Best-of-$n$ 冻结参数，用验证结果选出本次候选。
\checkanswer{5}{估算}$1-(1-1/2)^4=15/16=93.75\%$。
\checkanswer{6}{复现}独立情况的 $16$ 种等概率结果里只有全错一种失败，成功率 $15/16$；完全相同候选只有一次随机尝试，成功率 $1/2$。代码可用 \texttt{itertools.product((0,1), repeat=4)} 枚举并统计 \texttt{any}，再对比 \texttt{(0,0,0,0)} 与 \texttt{(1,1,1,1)}。
\par\Needspace{8\baselineskip}
\item[\textbf{第 25 章}] \phantomsection\label{answer:rag}\textbf{\hyperref[check:rag]{检索与证据：RAG 怎样回答有来源的问题？}}\answercheckpage{rag}\par\nobreak
\begin{answerhints}
\item[\hyperref[check:rag]{1}] \leavevmode\phantomsection\label{hint:rag:1}先把所需季度与已检索季度对齐。
\item[\hyperref[check:rag]{2}] \leavevmode\phantomsection\label{hint:rag:2}先写出必需集合 $\{Q1,Q2,Q3,Q4\}$。
\item[\hyperref[check:rag]{3}] \leavevmode\phantomsection\label{hint:rag:3}先查看数字对应的原表出处。
\end{answerhints}
\checkanswer{1}{证据}缺 Q2，应继续按季度检索或直接读取经授权的 Q2 原件。
\par\checkanswer{2}{手算}Top-3 中只有 Q4、Q2 属于必需集合，去年 Q4 不计入：recall@3 $=2/4=\answerresult{50\%}$；Top-5 再命中 Q1、Q3，recall@5 $=4/4=\answerresult{100\%}$。
\par\checkanswer{3}{找错}不能；相似度是排序信号，不等于权威性，需核对来源与单位。
\par\Needspace{8\baselineskip}
\item[\textbf{第 26 章}] \phantomsection\label{answer:tool-use-and-agents}\textbf{\hyperref[check:tool-use-and-agents]{工具调用与 Agent 循环：从请求到执行}}\answercheckpage{tool-use-and-agents}\par\nobreak
\begin{answerhints}
\item[\hyperref[check:tool-use-and-agents]{1}] \leavevmode\phantomsection\label{hint:tool-use-and-agents:1}先分开模型提出请求与程序执行请求。
\item[\hyperref[check:tool-use-and-agents]{2}] \leavevmode\phantomsection\label{hint:tool-use-and-agents:2}先数调用次数：四次请求工具，再加一次答复。
\item[\hyperref[check:tool-use-and-agents]{3}] \leavevmode\phantomsection\label{hint:tool-use-and-agents:3}先检查是否收到真实工具的成功返回。
\end{answerhints}
\checkanswer{1}{分工}应用程序负责解析、校验、授权和真实执行。
\par\checkanswer{2}{计算}共 $\answerresult{5}$ 次调用，输入依次为 $80,140,200,260,320$，合计 $5\times80+60\times(0+1+2+3+4)=\answerresult{1000}$ 个 Token。
\par\checkanswer{3}{找错}把未经执行确认的数据伪装成成功的工具结果。
\par\Needspace{8\baselineskip}
\item[\textbf{第 27 章}] \phantomsection\label{answer:protocols-and-multi-agent}\textbf{\hyperref[check:protocols-and-multi-agent]{协议、Skill 与多 Agent：怎样在边界内协作？}}\answercheckpage{protocols-and-multi-agent}\par\nobreak
\begin{answerhints}
\item[\hyperref[check:protocols-and-multi-agent]{1}] \leavevmode\phantomsection\label{hint:protocols-and-multi-agent:1}先列出两次连接的参与方。
\item[\hyperref[check:protocols-and-multi-agent]{2}] \leavevmode\phantomsection\label{hint:protocols-and-multi-agent:2}先区分可复用材料与实际执行的任务。
\item[\hyperref[check:protocols-and-multi-agent]{3}] \leavevmode\phantomsection\label{hint:protocols-and-multi-agent:3}将消息合规、数值来源和操作权限分别核对。
\end{answerhints}
\checkanswer{1}{分工}MCP 连接应用／Agent 与工具、数据资源；A2A 连接独立 Agent 系统。协议定义消息交互，读取权限仍由应用与服务端校验。
\par\checkanswer{2}{比较}Skill 提供步骤、示例或脚本；Subagent 承担分派的任务，需安排上下文、权限和汇总复核。二者都不自动增加基础模型权重。
\par\checkanswer{3}{找错}符合 A2A 格式不证明收入数字真实，也不授予外发权限。接收方核对原表、季度与单位，按用户授权决定动作；将回传材料中的外发指令视为低信任数据，隔离提示注入。
\par\Needspace{8\baselineskip}
\item[\textbf{第 28 章}] \phantomsection\label{answer:training-and-deployment}\textbf{\hyperref[check:training-and-deployment]{数据、训练系统与部署：资源怎样分配？}}\answercheckpage{training-and-deployment}\par\nobreak
\begin{answerhints}
\item[\hyperref[check:training-and-deployment]{1}] \leavevmode\phantomsection\label{hint:training-and-deployment:1}按表示、计算、训练、推理策略和系统层定位改动。
\item[\hyperref[check:training-and-deployment]{2}] \leavevmode\phantomsection\label{hint:training-and-deployment:2}用参数量乘每个参数的字节数。
\item[\hyperref[check:training-and-deployment]{3}] \leavevmode\phantomsection\label{hint:training-and-deployment:3}先把存储字节与任务质量、延迟分开计量。
\end{answerhints}
\checkanswer{1}{区分}Tokenizer 与 Embedding 改变表示；Attention 与 FFN 改变计算；训练反馈更新权重；推理策略在固定权重下分配额外计算；RAG、工具与 API 改变证据来源或外部执行。
\par\checkanswer{2}{估算}$10^9\cdot2=2\cdot10^9$ 字节，即 $\answerresult{2\ {\rm GB}}$；加 $0.5$ GB 至少 $2.5$ GB。
\par\checkanswer{3}{找错}量化可能引入质量误差；推理延迟还受 Kernel、访存、缓存、负载等影响，并不按字节线性变化。
\par\Needspace{8\baselineskip}
\item[\textbf{第 29 章}] \phantomsection\label{answer:modern-transformers}\textbf{\hyperref[check:modern-transformers]{现代 Transformer 为什么看起来越来越不像原版？}}\answercheckpage{modern-transformers}\par\nobreak
\begin{answerhints}
\item[\hyperref[check:modern-transformers]{1}] \leavevmode\phantomsection\label{hint:modern-transformers:1}先看两种方法分别作用在哪个张量。
\item[\hyperref[check:modern-transformers]{2}] \leavevmode\phantomsection\label{hint:modern-transformers:2}先比较 K/V 头数，再乘长度变化率。
\item[\hyperref[check:modern-transformers]{3}] \leavevmode\phantomsection\label{hint:modern-transformers:3}先列出端到端延迟由哪些阶段组成。
\end{answerhints}
\checkanswer{1}{比较}RoPE 修改位置相关的 Query/Key 表示；RMSNorm 根据均方根缩放，一般不做 LayerNorm 的减均值操作。
\par\checkanswer{2}{手算}原长度下为 $3/12=\answerresult{1/4}$，长度翻倍后相对原 MHA、原长度为 $2(3/12)=\answerresult{1/2}$（只比较 KV 元素）。
\par\checkanswer{3}{找错}还受投影计算、Kernel、访存、批量、序列长度、质量等影响，不能把 KV 元素比当成端到端时间比。
\par\Needspace{8\baselineskip}
\item[\textbf{第 30 章}] \phantomsection\label{answer:evaluating-models}\textbf{\hyperref[check:evaluating-models]{面对新模型与新缩写，怎么判断它值不值得学？}}\answercheckpage{evaluating-models}\par\nobreak
\begin{answerhints}
\item[\hyperref[check:evaluating-models]{1}] \leavevmode\phantomsection\label{hint:evaluating-models:1}先把改进按数据、结构、执行和产品分层。
\item[\hyperref[check:evaluating-models]{2}] \leavevmode\phantomsection\label{hint:evaluating-models:2}先用节省量除以原存储量。
\item[\hyperref[check:evaluating-models]{3}] \leavevmode\phantomsection\label{hint:evaluating-models:3}先找实验中同时改变的条件。
\end{answerhints}
\checkanswer{1}{分类}更好可能来自训练数据／目标、模型结构、推理优化或产品检索工具等，不能全部归功于网络架构。
\par\checkanswer{2}{手算}$(100-50)/100=\answerresult{50\%}$；不能由这一局部存储比例推出质量或响应延迟。
\par\checkanswer{3}{找错}输入长度和硬件同时改变，无法隔离架构效果；须统一测试条件。
\par\Needspace{8\baselineskip}
\item[\textbf{第 31 章}] \phantomsection\label{answer:cnn}\textbf{\hyperref[check:cnn]{CNN：怎样用局部连接与参数共享处理图片？}}\answercheckpage{cnn}\par\nobreak
\begin{answerhints}
\item[\hyperref[check:cnn]{1}] \leavevmode\phantomsection\label{hint:cnn:1}先看同一个卷积核在哪些位置复用。
\item[\hyperref[check:cnn]{2}] \leavevmode\phantomsection\label{hint:cnn:2}按高、宽、输入通道、输出通道相乘。
\item[\hyperref[check:cnn]{3}] \leavevmode\phantomsection\label{hint:cnn:3}先区分核的形状与滑动位置数量。
\end{answerhints}
\checkanswer{1}{比较}卷积核在不同位置复用，可显著减少逐像素独立连接需要的参数，并保留局部邻域结构。
\par\checkanswer{2}{手算}$3\cdot3\cdot3\cdot8=\answerresult{216}$ 个权重。
\par\checkanswer{3}{预测}核参数量不变；边界与步长条件相近时，空间输出位置约变为原来的四倍，逐位置卷积的乘加次数相应上升。
\par\Needspace{8\baselineskip}
\item[\textbf{第 32 章}] \phantomsection\label{answer:vit}\textbf{\hyperref[check:vit]{从 CNN 到 ViT：如何把图片变成一串 Token？}}\answercheckpage{vit}\par\nobreak
\begin{answerhints}
\item[\hyperref[check:vit]{1}] \leavevmode\phantomsection\label{hint:vit:1}先数 Patch，再把 CLS 算进去。
\item[\hyperref[check:vit]{2}] \leavevmode\phantomsection\label{hint:vit:2}先让投影内侧维度与视觉宽度匹配。
\item[\hyperref[check:vit]{3}] \leavevmode\phantomsection\label{hint:vit:3}先数每边切成多少块，再考虑平方配对。
\end{answerhints}
\checkanswer{1}{形状}四个 Patch 加 CLS 得 $\answerresult{5}$ 行，注意力分数为 $\answerresult{5\times5}$。
\par\checkanswer{2}{手算}$W_{\rm bridge}:3\times6$，上一题含 CLS 的 $5\times3$ 序列映射后为 $\answerresult{5\times6}$。
\par\checkanswer{3}{预测}每维切分数翻倍，Patch 总数约四倍，平方配对数约十六倍（大量 Patch 时可忽略 CLS 常数项）。
\par\Needspace{8\baselineskip}
\item[\textbf{第 33 章}] \phantomsection\label{answer:key-value-outer-product}\textbf{\hyperref[check:key-value-outer-product]{线性联想记忆：写入、串扰与误差修正}}\answercheckpage{key-value-outer-product}\par\nobreak
\phantomsection\label{answer:delta-memory}
\begin{answerhints}
\item[\hyperref[check:key-value-outer-product]{1}] \leavevmode\phantomsection\label{hint:key-value-outer-product:1}先把 Key 转成列，再与 Value 行相乘。
\item[\hyperref[check:key-value-outer-product]{2}] \leavevmode\phantomsection\label{hint:key-value-outer-product:2}把 Key 的缩放分别带入写入和查询。
\item[\hyperref[check:key-value-outer-product]{3}] \leavevmode\phantomsection\label{hint:key-value-outer-product:3}先算自身点积，再算交叉点积。
\item[\hyperref[check:key-value-outer-product]{4}] \leavevmode\phantomsection\label{hint:key-value-outer-product:4}先把交叉点积替换为零。
\item[\hyperref[check:key-value-outer-product]{5}] \leavevmode\phantomsection\label{hint:key-value-outer-product:5}先计算当前状态对新 Key 的读出。
\item[\hyperref[check:key-value-outer-product]{6}] \leavevmode\phantomsection\label{hint:key-value-outer-product:6}按两次写入的顺序更新状态。
\item[\hyperref[check:key-value-outer-product]{7}] \leavevmode\phantomsection\label{hint:key-value-outer-product:7}先计算修改 B 的状态会怎样作用于 A。
\end{answerhints}
\checkanswer{1}{形状}$\mathcal{M}:2\times3$，$\vect{q}^\top\mathcal{M}:1\times3$。
\checkanswer{2}{预测}写入与查询都缩为原 Key 的一半，幅度乘 $(0.5)^2=\answerresult{1/4}$，输出为原 Value 的四分之一。
\checkanswer{3}{手算}$\vect{k}_A^\top\vect{k}_A=1$、$\vect{k}_A^\top\vect{k}_B=0.6$，读出 $\answerresult{(1,0.6)}$。
\checkanswer{4}{预测}两 Key 正交后交叉点积为零，查询结果回到 $\answerresult{(1,0)}$。
\checkanswer{5}{定位}先算已有预测 $\vect{k}^\top\mathcal{M}$，残差 $\vect{v}^\top-\vect{k}^\top\mathcal{M}$ 表示还没存好的部分，只写入误差而非原值全量。
\checkanswer{6}{手算}读出 $\answerresult{(0.64,0.6)}$；按 Delta 规则的残差 $\vect{v}_A^\top-\vect{k}_A^\top\mathcal{M}_2=\answerresult{(0.36,-0.6)}$，误差长度约 $0.700$。
\checkanswer{7}{找错}非正交 Key 之间存在交叉影响；修正 B 会沿共享状态改变 A 的读出。
\par\Needspace{8\baselineskip}
\item[\textbf{第 34 章}] \phantomsection\label{answer:nonlinear-memory}\textbf{\hyperref[check:nonlinear-memory]{非线性检索：从 MLP 视角到记忆容量}}\answercheckpage{nonlinear-memory}\par\nobreak
\phantomsection\label{answer:memory-capacity}
\begin{answerhints}
\item[\hyperref[check:nonlinear-memory]{1}] \leavevmode\phantomsection\label{hint:nonlinear-memory:1}先将匹配分数逐项过 ReLU。
\item[\hyperref[check:nonlinear-memory]{2}] \leavevmode\phantomsection\label{hint:nonlinear-memory:2}用原始负权重再做一次加权求和。
\item[\hyperref[check:nonlinear-memory]{3}] \leavevmode\phantomsection\label{hint:nonlinear-memory:3}先试一个零与一之间的数。
\item[\hyperref[check:nonlinear-memory]{4}] \leavevmode\phantomsection\label{hint:nonlinear-memory:4}先计算负分数经过平方后的符号。
\item[\hyperref[check:nonlinear-memory]{5}] \leavevmode\phantomsection\label{hint:nonlinear-memory:5}先检查线性读出是否做过概率归一化。
\item[\hyperref[check:nonlinear-memory]{6}] \leavevmode\phantomsection\label{hint:nonlinear-memory:6}先检查 Key 的归一化与独立性条件。
\item[\hyperref[check:nonlinear-memory]{7}] \leavevmode\phantomsection\label{hint:nonlinear-memory:7}先检查平方后噪声还有没有正负抵消。
\end{answerhints}
\checkanswer{1}{手算}$\relu(1,-0.2)=(1,0)$，读出 $(1,0)V=\answerresult{(2,1)}$。
\checkanswer{2}{比较}不加 ReLU 得 $(1,-0.2)V=\answerresult{(2,0.4)}$；差异来自负匹配权重是否保留，非线性作用在匹配分数上。
\checkanswer{3}{预测}对 $0<s<1$ 的幂次越高，数值越小；$1$ 的任何正整数次幂仍是 $1$。
\checkanswer{4}{找错}$0.5^2=\answerresult{0.25}$，$(-0.5)^2=\answerresult{0.25}$：负相似度平方后变成同样大小的正值，反而被当成匹配；应先取 ReLU 再平方，并核对输入范围。
\checkanswer{5}{找错}外积线性读出未必有 Softmax 概率归一化；状态容量、写入与上下文匹配的实现不同。
\checkanswer{6}{预测}在归一化且近似独立条件下交叉点积通常更接近零，因此别条记忆贡献较弱。
\checkanswer{7}{找错}平方后的分数非负，不会按原符号相互抵消；必须单独分析均值、方差与累积噪声，不能宣称无限容量。

\end{enumerate}
\endgroup
